      第~(2) 条有五个条件。
      第一条是封闭性：若 \( c\in\Re \) 且 \( \vec{s}\in S \)，则
      \( c\cdot\vec{s}\in S \)，因为
      \( c\cdot\vec{s}=c\cdot\vec{s}+0\cdot\zero \)。
      第二，因为 \( S \) 中的运算继承自 \( V \)，所以对 \( c,d\in\Re \) 与 \( \vec{s}\in S \)，
      \( S \) 中的标量乘积 \( (c+d)\cdot\vec{s}\, \) 等于 \( V \) 中的乘积
      \( (c+d)\cdot\vec{s}\, \)，而它等于 \( V \) 中的
      \( c\cdot\vec{s}+d\cdot\vec{s}\, \)，进而等于 \( S \) 中的
      \( c\cdot\vec{s}+d\cdot\vec{s}\, \)。

      第三、第四和第五个条件的验证与刚才给出的第二个条件的验证类似。
    \end{answer}
  \item 
    证明每个向量空间只有一个平凡子空间。
    \begin{answer}
      前一小节的一道练习证明了每个向量空间只有一个零向量（也就是说，只有
      一个向量是该空间的加法单位元）。
      而平凡空间只有一个元素，且该元素必是这个（唯一的）零向量。
    \end{answer}
  \item 
    证明对向量空间的任意子集 \( S \)，张成的张成等于张成
    \( \spanof{ \spanof{S} }=\spanof{S} \)。
    （\textit{提示。}
    $\spanof{S}$ 的成员是 $S$ 的成员的线性组合。
    $\spanof{\spanof{S}}$ 的成员是 $S$ 的成员的线性组合的
    线性组合。）
    \begin{answer}
      如提示所言，根本的原因是第一章的线性组合引理。
      为给出完整证明，我们将证明这两个集合互相包含。

      第一个包含关系 \( \spanof{ \spanof{S} }\supseteq\spanof{S}  \)
      是一个更一般且显然的事实的特例：对向量空间的任意子集
      \( T \)，都有 \( \spanof{T}\supseteq T \)。

      对于另一个包含关系，
      即 \( \spanof{ \spanof{S} }\subseteq\spanof{S} \)，
      从 \( \spanof{S} \) 中取 $m$ 个向量，即
      \( c_{1,1}\vec{s}_{1,1}+\cdots+c_{1,n_1}\vec{s}_{1,n_1} \)、\ldots、
      \( c_{1,m}\vec{s}_{1,m}+\cdots+c_{1,n_m}\vec{s}_{1,n_m} \)，
      并注意它们的任意线性组合
      \begin{equation*}
         r_1(c_{1,1}\vec{s}_{1,1}+\cdots+c_{1,n_1}\vec{s}_{1,n_1})+\cdots
        +r_m(c_{1,m}\vec{s}_{1,m}+\cdots+c_{1,n_m}\vec{s}_{1,n_m})
      \end{equation*}
      是 \( S \) 中元素的线性组合
      \begin{equation*}
        = (r_1c_{1,1})\vec{s}_{1,1}+\cdots+(r_1c_{1,n_1})\vec{s}_{1,n_1}+\cdots
        +(r_mc_{1,m})\vec{s}_{1,m}+\cdots+(r_mc_{1,n_m})\vec{s}_{1,n_m}
      \end{equation*}
      因而在 \( \spanof{S} \) 中。
      也就是说，只需回想：线性组合的线性组合（$S$ 中成员的）
      仍是线性组合（同样还是 $S$ 中成员的）。
    \end{answer}
  \item 
    我们见过的所有子空间都在其描述中以某种方式用到了零。
    例如，\nearbyexample{ex:SubspacesRTwo} 中的子空间由 $\Re^2$ 中
    第二个分量为零的所有向量组成。
    相比之下，
    $\Re^2$ 中第二个分量为 1 的向量集合
    不构成子空间（它对标量乘法不封闭）。
    另一个例子是 \nearbyexample{ex:PlaneSubspRThree}，其中对向量的条件
    是三个分量之和为零。
    如果那里的条件
    改为三个分量之和为 1，
    那么它就不是子空间
    （同样，它不封闭）。
    不过，对零的依赖并非严格必需。
    考虑集合
    \begin{equation*}
       \set{\colvec{x \\ y \\ z}\suchthat x+y+z=1}
    \end{equation*}
    在这些运算下。
    \begin{equation*}
       \colvec{x_1 \\ y_1 \\ z_1}+\colvec{x_2 \\ y_2 \\ z_2}
       =\colvec{x_1+x_2-1 \\ y_1+y_2 \\ z_1+z_2}
       \qquad
       r\colvec{x \\ y \\ z}=\colvec{rx-r+1 \\ ry \\ rz}
    \end{equation*}
    \begin{exparts}
      \partsitem 证明它不是 $\Re^3$ 的子空间。
         （\textit{提示。}   参见 \nearbyexample{ex:OperNotInherit}）。
      \partsitem 证明它是一个向量空间。         
         注意由前一问，
         \nearbylemma{th:SubspIffClosed} 不能适用。  
      \partsitem 证明 $\Re^3$ 的任何子空间必定经过原点，
         因此 $\Re^3$ 的任何子空间在其描述中必定涉及零。
         逆命题成立吗？  
          $\Re^3$ 的任何包含原点的子集，
          在给定继承的运算后会成为子空间吗？
    \end{exparts}
    \begin{answer}
      \begin{exparts}
         \partsitem 它不是子空间，因为这些运算不是继承来的运算。  
           例如，在这个空间中，
           \begin{equation*}
             0\cdot\colvec{x \\ y \\ z}=\colvec[r]{1 \\ 0 \\ 0}
           \end{equation*}
           而这当然在 $\Re^3$ 中不成立。
         \partsitem 我们可以把证明对加法封闭的论证
           与证明对标量乘法封闭的论证合并成
           一个证明对两个向量的线性组合封闭的单一论证。
           若 $r_1,r_2,x_1,x_2,y_1,y_2,z_1,z_2$ 都属于 $\Re$，则      
           \begin{multline*}
              r_1\colvec{x_1 \\ y_1 \\ z_1}
              +r_2\colvec{x_2 \\ y_2 \\ z_2}
              =\colvec{r_1x_1-r_1+1 \\ r_1y_1 \\ r_1z_1}
               +\colvec{r_2x_2-r_2+1 \\ r_2y_2 \\ r_2z_2}                 \\
            =\colvec{r_1x_1-r_1+r_2x_2-r_2+1 \\ r_1y_1+r_2y_2 \\ r_1z_1+r_2z_2}
           \end{multline*} 
           （注意这个空间中加法的定义是：第一个分量按
           $(r_1x_1-r_1+1)+(r_2x_2-r_2+1)-1$ 的方式合并，
           所以最后一个向量的第一个分量并非表示
           `$\hbox{}+2$'）。
           将最后一个向量的三个分量相加得
           $r_1(x_1-1+y_1+z_1)+r_2(x_2-1+y_2+z_2)+1=r_1\cdot0+r_2\cdot0+1=1$。

           其余大部分条件的验证都很容易（尽管这些运算的古怪之处
           使它们无法按常规进行）。
           加法的交换律如下。
           \begin{equation*}
             \colvec{x_1 \\ y_1 \\ z_1}+\colvec{x_2 \\ y_2 \\ z_2}
             =\colvec{x_1+x_2-1 \\ y_1+y_2 \\ z_1+z_2}
             =\colvec{x_2+x_1-1 \\ y_2+y_1 \\ z_2+z_1}
             =\colvec{x_2 \\ y_2 \\ z_2}+\colvec{x_1 \\ y_1 \\ z_1}
           \end{equation*}
           加法的结合律则是
           \begin{equation*}
             (\colvec{x_1 \\ y_1 \\ z_1}+\colvec{x_2 \\ y_2 \\ z_2})
              +\colvec{x_3 \\ y_3 \\ z_3}
             =\colvec{(x_1+x_2-1)+x_3-1 \\ (y_1+y_2)+y_3 \\ (z_1+z_2)+z_3}
           \end{equation*}
           而
           \begin{equation*}
             \colvec{x_1 \\ y_1 \\ z_1}
             +(\colvec{x_2 \\ y_2 \\ z_2}+\colvec{x_3 \\ y_3 \\ z_3})
             =\colvec{x_1+(x_2+x_3-1)-1 \\ y_1+(y_2+y_3) \\ z_1+(z_2+z_3)}
           \end{equation*}
           且两者相等。
           关于这个加法运算的单位元 
           如下
           \begin{equation*}
             \colvec{x \\ y \\ z}+\colvec[r]{1 \\ 0 \\ 0}
             =\colvec{x+1-1 \\ y+0 \\ z+0}
             =\colvec{x \\ y \\ z}
           \end{equation*}
           加法逆元类似。
           \begin{equation*}
             \colvec{x \\ y \\ z}+\colvec{-x+2 \\ -y \\ -z}
             =\colvec{x+(-x+2)-1 \\ y-y \\ z-z}
             =\colvec[r]{1 \\ 0 \\ 0}
           \end{equation*}

           关于标量乘法的那些条件也容易验证。
           对第一个条件，
           \begin{equation*}
             (r+s)\colvec{x \\ y \\ z}
             =\colvec{(r+s)x-(r+s)+1 \\ (r+s)y \\ (r+s)z}
           \end{equation*}
           而
           \begin{multline*}
             r\colvec{x \\ y \\ z}+s\colvec{x \\ y \\ z}
             =\colvec{rx-r+1 \\ ry \\ rz}+\colvec{sx-s+1 \\ sy \\ sz}  \\
             =\colvec{(rx-r+1)+(sx-s+1)-1 \\ ry+sy \\ rz+sz}
           \end{multline*}
           两者相等。
           第二个条件比较
           \begin{equation*}
             r\cdot(\colvec{x_1 \\ y_1 \\ z_1}+\colvec{x_2 \\ y_2 \\ z_2})
             =r\cdot\colvec{x_1+x_2-1 \\ y_1+y_2 \\ z_1+z_2}
             =\colvec{r(x_1+x_2-1)-r+1 \\ r(y_1+y_2) \\ r(z_1+z_2)}
           \end{equation*}
           与
           \begin{multline*}
             r\colvec{x_1 \\ y_1 \\ z_1}+r\colvec{x_2 \\ y_2 \\ z_2}
             =\colvec{rx_1-r+1 \\ ry_1 \\ rz_1}
                     +\colvec{rx_2-r+1 \\ ry_2 \\ rz_2}                  \\
             =\colvec{(rx_1-r+1)+(rx_2-r+1)-1 \\ ry_1+ry_2 \\ rz_1+rz_2}
           \end{multline*}
           且它们相等。
           对第三个条件，
           \begin{equation*}
             (rs)\colvec{x \\ y \\ z}
             =\colvec{rsx-rs+1 \\ rsy \\ rsz}
           \end{equation*}
           而
           \begin{equation*}
             r(s\colvec{x \\ y \\ z})
             =r(\colvec{sx-s+1 \\ sy \\ sz})
             =\colvec{r(sx-s+1)-r+1 \\ rsy \\ rsz}
           \end{equation*}
           两者相等。
           对于被 $1$ 的标量乘法，我们有下式。
           \begin{equation*}
             1\cdot\colvec{x \\ y \\ z}
             =\colvec{1x-1+1 \\ 1y \\ 1z}
             =\colvec{x \\ y \\ z}
           \end{equation*}
           这样，这两个运算满足向量空间的所有条件。

           \textit{注记。}
           理解这个向量空间的一种方式是把它看作 $\Re^3$ 中的平面
           \begin{equation*}
             P=\set{\colvec{x \\ y \\ z}\suchthat x+y+z=0}
           \end{equation*}
           沿 $x$ 轴离开原点平移 $1$ 所得的平面。
           于是加法变成：~要加上这个空间的两个成员
           \begin{equation*}
             \colvec{x_1 \\ y_1 \\ z_1},\;\colvec{x_2 \\ y_2 \\ z_2}
           \end{equation*}
           （其中 $x_1+y_1+z_1=1$ 且 $x_2+y_2+z_2=1$）
           先把它们移回 $1$ 放入 $P$，再按通常方式相加，
           \begin{equation*}
             \colvec{x_1-1 \\ y_1 \\ z_1}+\colvec{x_2-1 \\ y_2 \\ z_2}
             =\colvec{x_1+x_2-2 \\ y_1+y_2 \\ z_1+z_2}
             \qquad\text{（在 $P$ 中）}
           \end{equation*}
           然后把结果沿 $x$ 轴再移出 $1$。
           \begin{equation*}
             \colvec{x_1+x_2-1 \\ y_1+y_2 \\ z_1+z_2}.
           \end{equation*}
           标量乘法类似。
         \partsitem 为使该子空间对继承的标量
           乘法封闭，设 $\vec{v}$ 是该子空间的一个成员，
           \begin{equation*}
             0\cdot\vec{v}=\colvec[r]{0 \\ 0 \\ 0}
           \end{equation*}
           它也必须是成员。

           逆命题不成立。
           下面是 $\Re^3$ 的一个包含原点的子集
           \begin{equation*}
             \set{\colvec[r]{0 \\ 0 \\ 0},\colvec[r]{1 \\ 0 \\ 0}}
           \end{equation*}
           （这个子集只有两个元素）但它不是子空间。
       \end{exparts}
     \end{answer}
  \item 
    我们可以为“零个向量之和等于零向量”这一约定给出一个理由。
    考虑三个向量的和 $\vec{v}_1+\vec{v}_2+\vec{v}_3$。
    \begin{exparts}
      \partsitem 这个三个向量的和与这三个中前两个的和之差是什么？
      \partsitem 前一个和与仅第一个向量的和之差是什么？
      \partsitem 前一个只含一个向量的和与零个向量的和之差应当是什么？
      \partsitem 那么零个向量的和应当定义为什么？
    \end{exparts}
    \begin{answer}
      \begin{exparts}
        \item $(\vec{v}_1+\vec{v}_2+\vec{v}_3)-(\vec{v}_1+\vec{v}_2)
                 =\vec{v}_3$
        \item $(\vec{v}_1+\vec{v}_2)-(\vec{v}_1)
                 =\vec{v}_2$
        \item 显然，$\vec{v}_1$。
        \item 取那个只含一项的和再相减，得
          （$\vec{v}_1)-\vec{v}_1=\zero$。
      \end{exparts}
    \end{answer}
  \item 
    一个空间由它的子空间决定吗？
    也就是说，如果两个向量空间有相同的子空间，这两个空间是否必定相等？
    \begin{answer}
      是的；任何空间都是它自身的子空间，所以每个空间都包含另一个。  
    \end{answer}
  \item 
     \begin{exparts}
      \partsitem 给出一个对标量乘法封闭
        但对加法不封闭的集合。
      \partsitem 给出一个对加法封闭但对标量乘法
        不封闭的集合。
      \partsitem 给出一个对两者都不封闭的集合。
    \end{exparts}
    \begin{answer}
      \begin{exparts}
        \partsitem \( \Re^2 \) 中 \( x \) 轴与 \( y \) 轴的并就是一例。
        \partsitem 整数集，作为 \( \Re^1 \) 的子集，就是一例。
        \partsitem \( \Re^2 \) 的子集 \( \set{\vec{v}} \) 就是一例，
          其中 $\vec{v}$ 是任意非零向量。
      \end{exparts}  
     \end{answer}
  \item 
    证明一组向量的张成不依赖于列举这些向量时的顺序，即
    该集合中向量被列举的顺序。
    \begin{answer}
      由于向量空间加法是交换的，重新排列被加项不会改变线性组合。  
    \end{answer}
  \item  
    空集的张成是哪个平凡子空间？
    是
    \begin{equation*}
      \set{\colvec[r]{0 \\ 0 \\ 0}}\subseteq \Re^3,
      \quad\text{或}\quad
      \set{0+0x}\subseteq \polyspace_1,
    \end{equation*}
    还是某个别的子空间？
    \begin{answer}
      我们总是在一个包围空间的背景下考虑这个张成。  
    \end{answer}
  \item   
    证明若一个向量在某个集合的张成中，则把该向量加进集合
    不会使张成变大。
    这在“仅当”方向上也成立吗？
    \begin{answer}
      “当”与“仅当”两个方向都成立。

      对于“当”的方向，
      设 \( S \) 是向量空间 \( V \) 的子集，并假设
      \( \vec{v}\in S \) 满足
      \( \vec{v}=c_1\vec{s}_1+\dots+c_n\vec{s}_n \)，其中
      \( c_1,\ldots,c_n \) 是标量，且
      \( \vec{s}_1,\ldots,\vec{s}_n\in S \)。
      我们必须证明 \( \spanof{S\union\set{\vec{v}} }=\spanof{S} \)。

      一侧的包含关系，
      \( \spanof{S}\subseteq\spanof{S\union\set{\vec{v}} } \) 是显然的。
      对于另一个方向，
      \( \spanof{S\union\set{\vec{v}} }\subseteq\spanof{S} \)，注意若一个
      向量属于左边的集合，则它形如
      \( d_0\vec{v}+d_1\vec{t}_1+\dots+d_m\vec{t}_m \)，其中诸 \( d \) 是
      标量，诸 \( \vec{t}\, \) 属于 \( S \)。
      把它改写为
      \( d_0(c_1\vec{s}_1+\dots+c_n\vec{s}_n)
      +d_1\vec{t}_1+\cdots+d_m\vec{t}_m \)，并注意
      该结果是 \( S \) 的张成中的一个成员。

      “仅当”显然成立\Dash 加入 \( \vec{v} \) 
      会把张成扩大到至少包含 \( \vec{v} \)。
    \end{answer}
  \recommended \item
    子空间是子集，所以我们自然会考虑“是……的子空间”这一关系
    与通常的集合运算如何相互作用。
    \begin{exparts}
      \partsitem 若 \( A,B \) 是向量空间的子空间，它们的交
        \( A\intersection B \) 必是子空间吗？
        总是？有时？从不？
      \partsitem 并 \( A\union B \) 必是子空间吗？
      \partsitem 若 \( A \) 是某个~$V$ 的子空间，它的补集 $V-A$
        必是子空间吗？
    \end{exparts}
    （\textit{提示。}   试用 \nearbyexample{ex:SubspRThree} 中
    的一些试验性子空间。）
    \begin{answer}
      \begin{exparts}
        \partsitem 总是。

          设 \( A,B \) 是 \( V \) 的子空间。
          注意 
          它们的交不空，因为两者都包含零向量。
          若 \( \vec{w},\vec{s}\in A\intersection B \) 且 \( r,s \) 是
          标量，则 \( r\vec{v}+s\vec{w}\in A \)，因为
          每个向量都属于 \( A \)，所以线性组合属于 \( A \)，
          出于同样的理由 \(r\vec{v}+s\vec{w}\in B \)。
          于是这个交是封闭的。
          这时 \nearbylemma{th:SubspIffClosed} 适用。
        \partsitem 有时（更准确地说，仅当 \( A\subseteq B \) 或
          \( B\subseteq A \)）。

          为看出答案不是“总是”，取 \( V \) 为 \( \Re^3 \)，
          取 \( A \) 为 $x$ 轴，\( B \) 为
          \( y \) 轴。
          注意
          \begin{equation*}
            \colvec[r]{1 \\ 0}\in A \text{ 且 }\colvec[r]{0 \\ 1}\in B
            \quad\text{而}\quad
            \colvec[r]{1 \\ 0}+\colvec[r]{0 \\ 1}\not\in A\union B
          \end{equation*}
          因为该和既不在 \( A \) 中也不在 \( B \) 中。

          答案不是“从不”，因为若 \( A\subseteq B \) 或
          \( B\subseteq A \)，则显然 \( A\union B \) 是子空间。

          为证明 \( A\union B \) 是子空间仅当一个子空间
          包含另一个，我们假设 \( A\not\subseteq B \)
          且 \( B\not\subseteq A \)，并证明 
          该并不是子空间。
          \( A \) 不是 \( B \) 的子集这一假设意味着存在 
          \( \vec{a}\in A \) 且 \( \vec{a}\not\in B \)。
          另一假设给出一个 \( \vec{b}\in B \)，满足
          \( \vec{b}\not\in A \)。
          考虑 \( \vec{a}+\vec{b} \)。
          注意该和不是 \( A \) 的元素，否则 
          \( (\vec{a}+\vec{b})-\vec{a} \) 将属于 \( A \)，但它不属于。
          类似地，该和也不是 \( B \) 的元素。
          因此该和不是 \( A\union B \) 的元素，所以该并
          不是子空间。
        \partsitem 从不。
          由于 \( A \) 是子空间，它包含零向量，因此
          作为 $A$ 的补集的集合 $V-A$ 不包含它。
          没有零向量，这个补集不可能是向量空间。
      \end{exparts}  
    \end{answer}
  \item 
    一个集合的张成依赖于包围它的空间吗？
    也就是说，若 \( W \) 是 \( V \) 的子空间，\( S \) 是
    \( W \) 的子集（因而也是 \( V \) 的子集），那么 \( S \) 在
    \( W \) 中的张成可能不同于 \( S \) 在 \( V \) 中的张成吗？
    \begin{answer}
      一个集合的张成不依赖于包围它的空间。
      \( S \) 中向量的线性组合给出同样的和，
      无论我们把运算视为 \( W \) 的还是 \( V \) 的，
      因为 \( W \) 的运算继承自 \(  V \)。  
    \end{answer}
  \item 
    “是……的子空间”这一关系是传递的吗？
    也就是说，若 $V$ 是 $W$ 的子空间，$W$ 是 $X$ 的
    子空间，$V$ 必是 $X$ 的子空间吗？ 
    \begin{answer}
      是传递的；
      应用 \nearbylemma{th:SubspIffClosed}。
      （你必须考虑下面这一点。
      设 \( B \) 是向量空间 \( V \) 的子空间，再设
      \( A\subseteq B\subseteq V \) 是一个子空间。
      \( A \) 从哪个空间继承它的运算？
      答案是这无关紧要\Dash 两种情形下 \( A \) 继承的
      都是相同的运算。）  
    \end{answer}
  \item
    由于“张成”是集合上的一种运算，我们自然会考虑
    它与通常的集合运算如何相互作用。
    \begin{exparts}
      \partsitem 若 \( S\subseteq T \) 是向量空间的子集，是否有
        \( \spanof{S}\subseteq\spanof{T} \)？
        总是？有时？从不？
      \partsitem 若 \( S,T \) 是向量空间的子集，是否有
        \( \spanof{S\union T}=\spanof{S}\union\spanof{T} \)？
      \partsitem 若 \( S,T \) 是向量空间的子集，是否有
        \( \spanof{S\intersection T}=\spanof{S}\intersection\spanof{T} \)？
      \partsitem 补集的张成等于张成的补吗？
    \end{exparts}
    \begin{answer}
      \begin{exparts}
         \partsitem 总是；
           若 \( S\subseteq T \)，则 \( S \) 中元素的线性组合
           也是 \( T \) 中元素的线性组合。
         \partsitem 有时（更准确地说，当且仅当 
           \( S\subseteq T \) 或 \( T\subseteq S \)）。

           答案不是“总是”，\( \Re^3 \) 中的这个例子
           表明了这一点
           \begin{equation*}
             S=\set{\colvec[r]{1 \\ 0 \\ 0},\colvec[r]{0 \\ 1 \\ 0}},\quad
             T=\set{\colvec[r]{1 \\ 0 \\ 0},\colvec[r]{0 \\ 0 \\ 1}}
           \end{equation*}
           因为有下式。
           \begin{equation*}
             \colvec[r]{1 \\ 1 \\ 1}\in\spanof{S\union T}
             \qquad
             \colvec[r]{1 \\ 1 \\ 1}\not\in\spanof{S}\union \spanof{T}
           \end{equation*}

           答案不是“从不”，因为若其中一个集合包含另一个，
           则等式显然成立。
           我们可以通过假设 \( S\not\subseteq T \)（意味着存在
           向量 \( \vec{s}\in S \) 满足 
           \( \vec{s}\not\in T \)）
           且 \( T\not\subseteq S \)（给出一个 \( \vec{t}\in T \)，满足
           \( \vec{t}\not\in S \)），并注意
           \( \vec{s}+\vec{t}\in\spanof{S\union T} \)，
           再证明 
           \( \vec{s}+\vec{t}\not\in\spanof{S}\union\spanof{T} \)，
           来刻画等式仅在其中一个集合包含另一个时才发生。
         \partsitem 有时。

           显然
           \( \spanof{S\intersection T}
             \subseteq\spanof{S}\intersection\spanof{T} \)
           因为来自 
           \( S\intersection T \)
           的向量的任何线性组合既是 \( S \) 中向量的组合，也是
           \( T \) 中向量的组合。

           反方向的包含并不总成立。
           例如，在 \( \Re^2 \) 中，取
           \begin{equation*}
             S=\set{\colvec[r]{1 \\ 0},\colvec[r]{0 \\ 1}},\quad
             T=\set{\colvec[r]{2 \\ 0}}
           \end{equation*}
           于是 \( \spanof{S}\intersection\spanof{T} \) 是 \( x \) 轴，
           但 \( \spanof{S\intersection T}  \) 是平凡子空间。

           要精确刻画等式何时成立并不容易。
           显然若其中一个集合包含另一个则等式成立，但由 \( \Re^3 \) 中的这个例子
           可知那不是“仅当”的。
           \begin{equation*}
             S=\set{\colvec[r]{1 \\ 0 \\ 0},\colvec[r]{0 \\ 1 \\ 0}},
             \quad
             T=\set{\colvec[r]{1 \\ 0 \\ 0},\colvec[r]{0 \\ 0 \\ 1}}
           \end{equation*}
         \partsitem 从不，因为补集的张成是子空间，而
           张成的补不是（它不包含零 
           向量）。
      \end{exparts}  
     \end{answer}
  % \item 
  %   Reprove \nearbylemma{le:SpanIsASubsp} without doing the
  %   empty set separately.
  %   \begin{answer}
  %     Call the subset \( S \).
  %     By \nearbylemma{th:SubspIffClosed},
  %     we need to check that 
  %     \( \spanof{S} \) is closed under linear combinations.
  %     If \( c_1\vec{s}_1+\dots+c_n\vec{s}_n,
  %       c_{n+1}\vec{s}_{n+1}+\dots+c_m\vec{s}_m\in\spanof{S} \) then
  %     for any \( p,r\in\Re \) we have
  %     \begin{multline*}
  %       p\cdot(c_1\vec{s}_1+\cdots+c_n\vec{s}_n)+
  %            r\cdot(c_{n+1}\vec{s}_{n+1}+\cdots+c_m\vec{s}_m)        \\
  %       =
  %       pc_1\vec{s}_1+\cdots+pc_n\vec{s}_n
  %         +rc_{n+1}\vec{s}_{n+1}+\cdots+rc_m\vec{s}_m
  %     \end{multline*}
  %     which is an element of \( \spanof{S} \).
  %     % (\textit{Remark.}
  %     % If the set $S$ is empty, then that
  %     % `if \ldots\ then \ldots' statement is vacuously true.)  
  %   \end{answer}
  \item 找一个对线性组合封闭 
    然而却不是向量空间的结构。
    % (\textit{Remark.} 
    % This is a bit of a trick question.)
    \begin{answer}
      要出现这种情况，保证加法与标量乘法运算合理性的条件中
      必须有一个被违反。
      考虑带这些运算的 \( \Re^2 \)。
      \begin{equation*}
        \colvec{x_1 \\ y_1}
        +\colvec{x_2 \\ y_2}
        =\colvec[r]{0 \\ 0}
        \qquad
        r\colvec{x \\ y}
        =\colvec[r]{0 \\ 0}
      \end{equation*}
      集合 $\Re^2$ 对这些运算封闭。 
      但它不是向量空间。
      \begin{equation*}
        1\cdot\colvec[r]{1 \\ 1}
        \neq\colvec[r]{1 \\ 1}
      \end{equation*}  
    \end{answer}
\index{subspace|)}
\index{vector space|)}
\end{exercises}
